\begin{table}[htbp]
\centering
\caption{Dose-Response: Does the Premium Rise with Nonlegal-Side Population?}
\label{tab:jrap_dose}
\small
\begin{tabular}{p{7.2cm}cccc}
\toprule
Specification & Coef. & SE (HC1) & SE (cl.) & $p$ (cl.) \\
\midrule
log(captive nonlegal population, 70mi) + full controls & $+0.035$ & -- & -- & $0.014$ \\
\quad excluding Sedgwick County, CO & $+0.027$ & -- & -- & $0.034$ \\
\bottomrule
\end{tabular}
\begin{minipage}{0.95\textwidth}
\vspace{0.3em}
\footnotesize
\textit{Notes:} Outcome is log(revenue per active dispensary). The captive-population measure sums the
population of \emph{all} nonlegal-state counties within 70 miles of a given county (not merely the nearest
one), via a radius query against national county population centroids, capturing $P_i$ from the conceptual
framework directly. Only 10 of the 71 Colorado and Washington counties in this panel have any nonlegal
population within 70 miles at all, so this reweights the same identifying set of border counties by market
size rather than adding new geographic variation; it is nonetheless the most robust mechanism evidence in
this paper, since it survives county-clustered inference where a simpler continuous-distance measure
(Table~\ref{tab:jrap_border}) does not.
\end{minipage}
\end{table}
